% §9.6 Bibliographic Notes
%
% Cross-reference resolution status:
%
% ALREADY RESOLVED (within chapter 9, drafted in Tasks 2-6 of Plan 9):
%   \Cref{sec:atom-to-function}                  -- chapter 9 §9.1
%   \Cref{sec:n-ary-channels}                    -- chapter 9 §9.2
%   \Cref{def:n-ary-channel}                     -- chapter 9 §9.2 definition
%   \Cref{sec:n-ary-kronecker}                   -- chapter 9 §9.3
%   \Cref{thm:n-ary-kronecker}                   -- chapter 9 §9.3 theorem
%   \Cref{cor:n-ary-parameter-count}             -- chapter 9 §9.3 corollary
%   \Cref{sec:n-ary-capacity}                    -- chapter 9 §9.4
%   \Cref{prop:n-ary-capacity}                   -- chapter 9 §9.4 proposition
%   \Cref{sec:bridge-to-sketches}                -- chapter 9 §9.5
%
% ALREADY RESOLVED (chapters 4, 8):
%   \Cref{sec:kronecker}                         -- chapter 4 §4.4
%   \Cref{thm:kronecker}                         -- chapter 4 §4.4 theorem
%   \Cref{sec:bsc-capacity}                      -- chapter 8 §8.2
%   \Cref{sec:entropy-refresher}                 -- chapter 8 §8.1
%
% FORWARD REFERENCES (resolve in later plans):
%   \Cref{ch:sketches}                           -- chapter 10 (Plan 0 stub)
%
% Citations resolved by new references.bib entries added in this task:
%   \cite{towell-bernoulli-maps}                 -- Towell approximate-maps manuscript
%
% Citations resolved by prior references.bib entries:
%   \cite{cover-thomas}                          -- already in references.bib (Plan 0)
%   \cite{shannon-1948}                          -- already in references.bib (Plan 0)

\section*{Bibliographic Notes}
\addcontentsline{toc}{section}{Bibliographic Notes}
\label{sec:bib-notes-ch9}

\paragraph{The approximate-maps manuscript.}
The reformulation of the noisy-bit atom $\Bernoulli{\BoolType}$ into a
general type-lifted noisy value $\Bernoulli{T}$, together with the
n-ary channel matrix $Q$ that organizes its per-block conditional
output distributions, follows \textcite{towell-bernoulli-maps}. That
manuscript is the primary source for the chapter: the move from the
predicate framing of Part~II to the $T$-valued framing of Part~III in
\Cref{sec:atom-to-function}, the formal $K \times L$ row-stochastic
channel of \Cref{def:n-ary-channel}, the n-ary Kronecker factorization
of \Cref{thm:n-ary-kronecker} and its parameter-count corollary
\Cref{cor:n-ary-parameter-count}, and the capacity formulation of
\Cref{sec:n-ary-capacity} are all taken from it. The chapter's
contribution beyond the manuscript is the pedagogical scaffolding: the
explicit recall of \Cref{sec:kronecker} and \Cref{sec:bsc-capacity} as
the model-order-2 special cases that the chapter generalizes, the
3-class classifier and integer-valued counter as concrete instances of
$\Bernoulli{T}$ for non-Boolean $T$, and the bridge-to-sketches
discussion of \Cref{sec:bridge-to-sketches} that prepares
chapter~\ref{ch:sketches} by exhibiting CountMin, HyperLogLog, and
MinHash as $\Bernoulli{T}$ instances of the same shape.

\paragraph{N-ary channels and the Blahut-Arimoto algorithm.}
The information-theoretic background for the n-ary discrete memoryless
channel is standard textbook material, and the depth reference for
readers who want more than the working summary of
\Cref{sec:n-ary-capacity} is \textcite{cover-thomas}. Their chapter~7
develops the discrete memoryless channel, defines capacity as the
supremum of mutual information over input distributions on a finite
alphabet, proves the noisy-channel coding theorem in both directions,
and treats the symmetric and weakly symmetric special cases whose
closed-form capacities the chapter quotes. Their chapter~10 covers
the Blahut-Arimoto algorithm, the iterative coordinate ascent on the
KL-divergence form of mutual information that the chapter's notebook
implements, with proofs of monotone improvement and convergence to the
capacity-achieving input distribution. The Blahut-Arimoto recipe is
the default numerical tool for any channel that does not admit a
closed-form capacity; it is what makes the general n-ary capacity of
\Cref{prop:n-ary-capacity} a computable quantity rather than a
definition.

\paragraph{Shannon's original framework.}
The channel-capacity vocabulary that chapter~8 \S8.2 introduces for
the binary symmetric channel and that this chapter generalizes to
arbitrary $K \times L$ matrices originates in
\textcite{shannon-1948}. The 1948 paper introduces the entropy
$\Entropy(X)$ recalled in \Cref{sec:entropy-refresher}, the channel
capacity $C = \max_{P_X} \MutualInfo(X; Y)$ that
\Cref{sec:n-ary-capacity} specializes to the n-ary discrete case, and
the source-coding theorem that connects the two through the operational
characterization of capacity as the largest rate at which information
can be transmitted with vanishing error. Shannon's framing of the
channel as a probabilistic mapping from input alphabet to output
alphabet, conditional on the input distribution and producing an
output distribution governed by a fixed transition kernel, is exactly
the framing that \Cref{def:n-ary-channel} adopts: the chapter's $Q$ is
Shannon's transition matrix, the chapter's $K$ and $L$ are the
cardinalities of his input and output alphabets, and the chapter's
capacity is his capacity. The 1948 paper remains worth reading on its
own terms, both for the conceptual clarity of its development and for
the perspective it gives on how thoroughly Shannon's vocabulary has
shaped every subsequent treatment of noisy information transmission,
this chapter included.

\paragraph{Forward to chapter 10.}
The sketch literature that \Cref{sec:bridge-to-sketches} previews
(CountMin, HyperLogLog, MinHash, and the streaming-algorithms
tradition that produced them) is taken up in chapter~\ref{ch:sketches}
\S10.6, where the original papers are cited alongside the survey
treatments that have become the standard secondary sources. The
strategy of chapter~9 is deliberately to fix the framework first and
defer the named constructions: the three sketches are presented in
chapter~10 as the first non-Boolean instances of $\Bernoulli{T}$ that
exercise the full $K \times L$ channel machinery developed here, and
the bibliographic discussion belongs with the constructions rather
than with the abstract framework that subsumes them.
